\documentclass[12pt, a4paper,answers]{exam}
\usepackage{graphicx,fullpage,complexity}
\usepackage{amsmath,amsthm,amssymb}
\usepackage{xcolor}
\definecolor{iitpkdcolor}{cmyk}{0,0.42,1,0} %iitpkd logo color
\usepackage[normalem]{ulem}
\usepackage{import}

\renewcommand\ULthickness{2.0pt}   %%---> For changing thickness of underline
\setlength\ULdepth{1.3ex}%\maxdimen ---> For changing depth of underline
\newcommand{\N}{\mathbb{N}}
\newcommand{\HP}{\mathsf{HP}}
\newcommand{\ALL}{\mathsf{ALL}}
\newcommand{\TOTAL}{\mathsf{TOTAL}}
\renewcommand{\D}{\mathcal{D}}
\newcommand{\SD}{\mathcal{SD}}
%\newcommand{\BP}{mathsf{BP}}
%\newcommand{\}{def}
\begin{document}
	\newcommand{\lem}{\le_m}
\newcommand{\imp}{$\Rightarrow$}
\newcommand{\dimp }{$\Leftrightarrow$}
\newcommand{\la }{\langle}
\newcommand{\ra }{\rangle}
\newcommand{\calp}{\mathcal{P}}
\newcommand{\calpc}{\overline{\mathcal{P}}}

	%\thispagestyle{empty}
	\noindent
	\begin{minipage}[l]{0.1\textwidth}
		\noindent
		\includegraphics[width=2.4\textwidth]{iitpkd}
	\end{minipage}
	\hfill
	\begin{minipage}[c]{0.8\textwidth}
		\begin{center}
			{\large \textsc{\textcolor{iitpkdcolor}{Indian Institute of Technology Palakkad}} \par
				\small	Department of Computer Science and Engineering	\par
				\large	CS5616 Computational Complexity \par 
				\small January -- May 2024}
		\end{center}
	\end{minipage}
	\par
\vspace{2mm}
\hrule
\vspace{2mm}
\begin{minipage}{0.9\textwidth} 
	\textsf{Problem Set} -- 3  \hfill  \textsf{Total Points} -- 50
	
	{\small \textsf{Name}: Neeraj Krishna N   \hfill \small \textsf{Given on} 29 Feb}

	
	{\small \textsf{Roll no}: 112101033 \hfill \small \textsf{Due on} 11 Mar}
\end{minipage}
	\vspace{0.2in}
\noindent

\textbf{Instructions}
\begin{itemize}   \setlength\itemsep{0.1mm}
	\item {\sf  Use of resources other than class notes and references is forbidden.}
	\item {\sf Collaboration is not allowed. Credit will be given for attempts and partial answers.}
\end{itemize}


\begin{questions}

\question[10] (\textbf{RE vs co-RE}) Show that the set \[ \{M \mid M \text{ halts on all inputs of length less than $42$}\}\] is recursively enumerable, but is its complement is not.

\begin{solution}
	\import*{./q1}{q1.tex}
\end{solution}

	\question[10] (\textbf{Alternate definition for $\Delta_i$}) Let $A$ be any language. Define $\D^A$ be the class of all languages $L$ such that $L$ is decidable in $A$. Similarly, $\SD^A$ be the class of all $L$ such that $L$ is semi-decidable in $A$ and $\co\SD^A $ be the class of all languages whose complement is in $\SD^A$.
	
	\begin{parts}
	\part[5] Show that $\D^A = \SD^A \cap \co\SD^A$.
	
	\part[5] For any $i \ge 1$, by definition, $\Delta_i = \Sigma_i \cap \Pi_i$. Show that  \[ \Delta_i = \{L \mid \text{ there exists } A \in \Sigma_{i-1} \text{ such that } L \text{ is decidable in A}\}. \]
	\end{parts}

\begin{solution}
	\begin{enumerate}
		\item[(a)] Non-deterministic Turing Machine is a $9$-tuple $M = (Q, \Sigma, \Gamma, \vdash, \sqcup, \delta, s, t, r)$ where

\begin{itemize}
    \item $Q$ is a finite set of states
    \item $\Sigma$ is a finite set of input alphabets
    \item $\Gamma$ is the finite set of tape symbols
    \item $\vdash \in \Gamma \setminus \Sigma$ is the left end marker
    \item $\sqcup \in \Gamma \setminus \Sigma$ is the blank symbol
    \item $\delta : Q \times \Gamma  \rightarrow 2 ^ {Q \times \Gamma \times \{L, R\}}$ is the transition relation
    \item $s \in Q$ is the start state
    \item $t \in Q$ is the accept state
    \item $r \in Q$ is the reject state
\end{itemize}

\textbf{Configuration} :
\newline
A configuration $c$ is a tuple of $(q, w\sqcup ^ \omega, n)$ where $q \in Q$, $w \in \Gamma^*$, $n \in \N$

\textbf{Next Configuration Relation} :
\newline
It is the single step transition from one configuration to a set of configurations.
$(p, z, n) \rightarrow (q, z', n')$ where 
$ n' = n + 1$ if head moved right, else if head moved left, $n' = n - 1$, else $n' = n$ and $((p, z), (q, z')) \in \delta$

\textbf{Acceptance} : 
\newline
Turing machine M, is said to accept an input $x$, if after a finite number of steps, the configuration becomes $(t, y, n)$ for some $y \in \Gamma ^ *$ and $n \in \N$



		\item[(b)] 


Let $L_1 = 0^*$ and $L_2 = 1^*0^*$

Here, $L_1 \subseteq L_2$

The reduction machine on input $x$, an instance of $L_2$ produces an instance $y$ of $L_1$

To prove : $L_1 \lem L_2$

Let $M_2$ be the total turing machine which accepts $L_2$, we know it exists, because $L_2$ is \textbf{regular}

The reduction machine $M_\sigma$ Is

``

On input $x$
\begin{enumerate}
    \item Run $M_2$ on $x$
    \item If $M_2$ accepts, output $0$
    \item If $M_2$ rejects, output $1$
\end{enumerate}

''

Now, we have to prove
$x \in L_2 \Leftrightarrow M_\sigma(x) \in L_1$

$x \in L_2$

\imp
$M_2$ accepts $x$

\imp
$M_\sigma$ outputs $0$ which is in $L_1$

\imp
$0 \in L_1$

\imp
$M_\sigma(x) \in L_1$


$x \notin L_2$

\imp
$M_2$ rejects $x$, since $M_2$ is total

\imp
$M_\sigma(x) = 1$ which is not in $L_1$

\imp
$M_\sigma(x) \notin L_1$


Hence given $L_1 = 0^*$ and $L_2 = 1^*0^*$ such that $L_1 \subseteq L_2$ and $L_2 \lem L_1$

















	\end{enumerate}
\end{solution}

% Uncomment the following to answer
%\begin{solution}
%	Write your answer here.
%	\begin{description}
%		\item[(a)] Answer to first part
%		\item[(b)] Answer to second part
%	\end{description}
%\end{solution}	

\question[10] (\textbf{Closure properties of $\Sigma_n, \Pi_n$}). Fix any $i \ge 1$. Show that $\Sigma_i$ as well as $\Pi_i$ are closed under intersection and union.
	%%	 %Uncomment the following to answer
	%	\begin{solution}
		%		Write your answer here.
		%	\end{solution}
		
	
	\begin{solution}
		
\textbf{Claim 1} : $\s{i}$ closed under union \imp
$\pii{i}$ closed under intersection

\textbf{Claim 2} : $\s{i}$ closed under intersection \imp
$\pii{i}$ closed under union

Proof of Claim 1 : 


$\s{i}$ closed under union

\imp
$\forall L_1, L_2 \in \s{i}, L_1 \cup L_2 \in \s{i}$

\imp
$\forall L_1, L_2 \in \s{i}, \ovl{(\ovl{L_1} \cap \ovl{L_2})} \in \s{i}$, by De Morgan's Law

\imp
$\forall \ovl{L_1}, \ovl{L_2} \in \pii{i}, {(\ovl{L_1} \cap \ovl{L_2})} \in \pii{i}$

\imp
$\forall L_1', L_2' \in \pii{i}, {(L_1' \cap L_2')} \in \pii{i}$

\imp
$\pii{i}$ closed under intersection

Proof of Claim 2 : 


$\s{i}$ closed under intersection

\imp
$\forall L_1, L_2 \in \s{i}, L_1 \cap L_2 \in \s{i}$

\imp
$\forall L_1, L_2 \in \s{i}, \ovl{(\ovl{L_1} \cup \ovl{L_2})} \in \s{i}$, by De Morgan's Law

\imp
$\forall \ovl{L_1}, \ovl{L_2} \in \pii{i}, {(\ovl{L_1} \cup \ovl{L_2})} \in \pii{i}$

\imp
$\forall L_1', L_2' \in \pii{i}, {(L_1' \cup L_2')} \in \pii{i}$

\imp
$\pii{i}$ closed under union

Suffices to prove

$\s{i}$ is closed under intersection and union

We'll argue using induction

\textbf{Base case} : $i = 0$
$\s{0} = \Delta_0 = \Delta_1$ which is the set of decidable languages and we already know that decidable languages are closed under union and intersection

\textbf{Strong Induction Hypothesis}
The below statement is true for all $i \leq n$
$$
\s{i} \text{ is closed under intersection and union}
$$
The below statement also follows from Claim 1 and claim 2
$$
\pii{i} \text{ is closed under intersection and union}
$$

Therefore we can consider both statements as Strong induction Hypothesis

\textbf{To prove for $i = n + 1$}

Let $L_1, L_2 \in \s{n + 1}$

\imp
$\exists A \in \pii{n}, A \subseteq \Sigma^* \times \Sigma^*, x \in L_1 \Leftrightarrow \exists y_1, (x, y_1) \in A$ and

$\exists B \in \pii{n}, B \subseteq \Sigma^* \times \Sigma^*, x \in L_2 \Leftrightarrow \exists y_2, (x, y_2) \in B$

% \imp
% $\exists N^{\ovl{A}}$ accepting $L_1$ and 
% $\exists N^{\ovl{B}}$ accepting $L_2$

\imp
$\exists (A \cup B) \in \pii{n}$ (by induction hypothesis) such that 

$
x \in L_1 \cup L_2 \Leftrightarrow \exists y, (x, y) \in (A \cup B)$

\imp
$L_1 \cup L_2 \in \pii{n + 1}$

and thus proved $\s{n + 1}$ is closed under union which proves $\pii{n + 1}$ is closed under intersection from the claims

Let $L_1, L_2 \in \s{n + 1}$

\imp
$\exists A \in \pii{n}, A \subseteq \Sigma^* \times \Sigma^*, x \in L_1 \Leftrightarrow \exists y_1, (x, y_1) \in A$ and

$\exists B \in \pii{n}, B \subseteq \Sigma^* \times \Sigma^*, x \in L_2 \Leftrightarrow \exists y_2, (x, y_2) \in B$

% \imp
% $\exists N^{\ovl{A}}$ accepting $L_1$ and 
% $\exists N^{\ovl{B}}$ accepting $L_2$

\imp
$\exists (A \cap B) \in \pii{n}$ (by induction hypothesis) such that 

$
x \in L_1 \cap L_2 \Leftrightarrow \exists y, (x, y) \in (A \cap B)$

\imp
$L_1 \cap L_2 \in \pii{n + 1}$


and thus proved $\s{n + 1}$ is closed under intersection which proves $\pii{n + 1}$ is closed under union from the claims

therefore we have proved that if the statement holds for all $i \leq n$, it holds for $i = n + 1$

and thus by principle of mathematical induction it is proved that 

$$
\forall i \in \N, \s{i} \text{ and } \pii{i} \text{ are closed under union and intersection}
$$



















	\end{solution}
	
\question [20] (\textbf{Rice's theorem}) Identify if the following are (0) properties of SD languages, (1) non-trivial properties and (2) non-monotone properties. If (1) / (2) is true, apply Rice's theorems suitably and give your conclusions. A direct use of diagonalisation or reductions does not fetch any credit.

\begin{parts}
	\part[4] given a Turing machine $M$, $L(M)$ is not regular ?
	\part[4] given a Turing machine $M$, does there exist a non-empty regular set $L'$ such that $L' \subseteq L(M)$ ?
	\part[4] given $M$, does $M$ represent a DFA that accepts some string with equal number of $0$s and $1$s ?
	\part[4] given a Turing machine $M$, is $L(M) \in \Pi_{42}$ ?
\end{parts}

% Uncomment the following to answer
%\begin{solution}
%	Write your answer here.
%	\begin{description}
	%		\item[(a)] Answer to first part
	%		\item[(b)] Answer to second part
	%	\end{description}
%\end{solution}	

\begin{solution}
	\begin{enumerate}
		\item[(a)] 
$
L = \{
    M \mid L(M) \text{ is not regular}
\}
$

is a property of SD language

If there are 2 machines $M_1$ and $M_2$ accepting the same language, then it cannot be the case that $L(M_1)$ is regular and $L(M_2)$ is not regular

This is non trivial since there are turing machines which accept regular languages and turing machines which accept non-regular languages

This is also non-monotone since, there exists $L_1, L_2$ such that $L_1 \subseteq L_2$ and $\calp(L_1) = 1$ and 
$\calp(L_2) = 0$

$L_2 = \Sigma^* $ and $L_1 = \{
    0^n1^n \mid n \geq 0
\}$

here $L_1$ is not regular which satisfies the property whereas $L_2$ is regular and a strict superset of $L_1$ which does not satisfy the property

Since the property is non-trivial, by Rice's Theorem 1, $L$ is not decidable

Since the property is non-monotone, by Rice's Theorem 2,
$L$ is not semi-decidable












		\item[(b)] ??????
\newline
This is undecidable because, we can show that if the given problem is decidable, then $HP$ is decidable.
Let us take a turing machine $M_\sigma$ which on input $\langle M, x \rangle$, provides a description of another turing machine $M'$ which is 
``
''





		\item[(c)] The language is 

























		\item[(d)] 
The language is 

$
L = \{
    M \mid L(M) \in \pii{42}
\}
$

This is a property of language since it cannot be the case that there exists $M_1, M_2$ accepting the same language but $L(M_1) \in \pii{42}$ and $L(M_2) \notin \pii{42}$ or vice-versa

This is non-trivial because there exists turing machines accepting $L \in \pii{42}$ and turing machines accepting languages $L \notin \pii{42}$

Hence this is undecidable by Rice's Theorem 1

This is also non-monotone because, let $L_1 \in \pii{42}$ and 
$L_2 \in \pii{43} \setminus \pii{42}$ 

$\calp(L_1) = 1$ and $\calp(L_1 \cup L_2) = 0$ since union of $L_1$ and $L_2$ cannot be in $\pii{42}$ since $L_2$ is in $\pii{43} \setminus \pii{42}$


















	\end{enumerate}
\end{solution}


\end{questions}

\end{document}
